% Worked masking example: present_sbox MIG
% Fault site g = new_n12 (MAJ3 of pi1, pi2, new_n11)
% Downstream gate h = new_n15 (MAJ3 of new_n10, ~new_n12, new_n14)
% g has exactly one consumer (h), so dp(g) is determined entirely by h's masking.

\begin{figure*}[tbp]
  \centering
  \begin{tikzpicture}[
    >=stealth,
    every node/.style={font=\small},
    gate/.style={draw, rounded corners=1pt, minimum width=1.6cm,
                 minimum height=0.7cm, fill=blue!8},
    fgate/.style={draw, rounded corners=1pt, minimum width=1.6cm,
                  minimum height=0.7cm, fill=red!10, draw=red!70},
    wire/.style={->, thick},
    inv/.style={circle, inner sep=0pt, minimum size=4pt, draw, fill=white},
  ]
    % Primary inputs feeding g
    \node (pi1)   at (-3.2,   1.0) {\scriptsize \texttt{pi1}};
    \node (pi2)   at (-3.2,   0.3) {\scriptsize \texttt{pi2}};
    \node (n11)   at (-3.2,  -0.4) {\scriptsize \texttt{new\_n11}};

    % Fault gate g
    \node[fgate] (g) at (-1.0, 0.3) {\textsc{maj3} $g$};
    \node[above=2pt of g, font=\scriptsize\itshape] {fault site};

    \draw[wire] (pi1) -- (g.west |- pi1);
    \draw[wire] (pi2) -- (g.west);
    \draw[wire] (n11) -- (g.west |- n11);

    % Other inputs to h
    \node (n10)   at (1.5,  1.0) {\scriptsize \texttt{new\_n10}};
    \node (n14)   at (1.5, -0.4) {\scriptsize \texttt{new\_n14}};

    % Downstream gate h (sole consumer of g)
    \node[gate] (h) at (3.8, 0.3) {\textsc{maj3} $h$};

    % Inverter bubble on g -> h edge (h treats g's contribution complemented)
    \node[inv] (bubble) at (2.7, 0.3) {};

    \draw[wire] (g.east) -- (bubble.west);
    \draw[wire] (bubble.east) -- (h.west);
    \draw[wire] (n10) -- (h.west |- n10);
    \draw[wire] (n14) -- (h.west |- n14);

    % Output
    \node (out) at (6.2, 0.3) {\scriptsize \texttt{new\_n15}};
    \draw[wire] (h.east) -- (out.west);

    % Annotation: masking rule
    \node[align=center, font=\scriptsize, draw=gray!50,
          rounded corners=2pt, fill=gray!5, inner sep=4pt]
      (note) at (2.4, -1.7)
      {When \texttt{new\_n10} $=$ \texttt{new\_n14} (agree): non-controlling pattern, fault masked.\\[-1pt]
       When \texttt{new\_n10} $\neq$ \texttt{new\_n14} (disagree): fault propagates. (8\,/\,16 each.)};
    \draw[densely dotted, gray!70] (h) -- (note);
  \end{tikzpicture}
  \caption{Worked masking example from the \texttt{present\_sbox} MIG.
    Gate $g$ (fault site) feeds a single downstream MAJ3 gate $h$
    through a complemented edge. A bit-flip at $g$ propagates to $h$'s
    output iff the other two inputs of $h$ \emph{disagree}; when they
    \emph{agree}, they form the non-controlling pattern and mask the
    fault. Over the 16 exhaustive input vectors, the other two inputs
    agree on 8 vectors (fault masked) and disagree on 8 (fault propagates). Because $g$ has $h$ as its sole consumer, this
    completely determines $g$'s observability: $\mathrm{dp}(g) = 8/16
    = 0.5$. The per-vector trace in Table~\ref{tab:masking-trace}
    enumerates all 16 vectors.}
  \label{fig:masking-example}
\end{figure*}
